\setcounter{section}{0}

\Needspace{5\baselineskip}
\hypertarget{architectural-approaches-to-embodied-ai-models}{%
\section{Architectural Approaches to Embodied AI Models}\label{architectural-approaches-to-embodied-ai-models}}

\Needspace{5\baselineskip}
\hypertarget{i.-model-architecture-approaches}{%
\subsection{I. Model Architecture Approaches}\label{i.-model-architecture-approaches}}

\Needspace{5\baselineskip}
\hypertarget{traditional-modular-approach}{%
\subsubsection{(1) Traditional Modular Approach}\label{traditional-modular-approach}}

Core logic: The classical decoupled ``perception--decision--control'' paradigm.

Workflow: The vision/sensor module handles object detection and mapping (such as YOLO and SLAM), the decision module handles global and local path planning, and traditional low-level control algorithms (such as MPC and PID) drive the motors to execute actions.

Characteristics: Highly interpretable and extremely mature in engineering deployment, but fundamentally lacking intelligence. Substantial information is lost across modules, making it difficult to handle long-tail, unstructured scenarios or develop general-purpose generalization capabilities.

\Needspace{5\baselineskip}
\hypertarget{hierarchical-model-approach}{%
\subsubsection{(2) Hierarchical Model Approach}\label{hierarchical-model-approach}}

Core logic: Introduce heterogeneous collaboration between a ``brain'' and a ``cerebellum.''

Workflow: A top-level vision--language model (VLM/LLM) acts as the ``brain,'' receiving multimodal environmental inputs and human instructions, performing high-dimensional semantic reasoning and task decomposition, and outputting intermediate representations (such as Python code, waypoint positions, or discrete skill/primitive instructions). A low-level control policy acts as the ``cerebellum,'' converting these representations into actions.

Characteristics: Balances LLMs' generalizable reasoning with the low-level execution precision of RL/control algorithms. However, it is not end-to-end, is difficult to optimize in a closed loop, and is limited in its level of intelligence.

\Needspace{5\baselineskip}
\hypertarget{visionlanguageaction-vla-approach}{%
\subsubsection{(3) Vision--Language--Action (VLA) Approach}\label{visionlanguageaction-vla-approach}}

Core logic: Establish an end-to-end mapping from multimodal sensor inputs to robot actions. Vision--language--action (VLA) models are the mainstream of this approach.

Workflow: Images, text instructions, and proprioceptive states (such as current joint angles) are uniformly converted into tokens and fed into a Transformer-based autoregressive or diffusion model, directly producing continuous or discrete action tokens with multiple degrees of freedom.

Characteristics: This is currently a relatively mainstream method. Because data is limited, most current VLAs are obtained by fine-tuning VLMs. Their language pretraining allows them to capture high-level semantics, but their fine-grained predictive capability is weaker. VLAs also still face insufficient generalization when the physical environment changes and transferability issues across different hardware.

\Needspace{5\baselineskip}
\hypertarget{worldaction-model-wam-approach}{%
\subsubsection{(4) World--Action Model (WAM) Approach}\label{worldaction-model-wam-approach}}

Core logic: Combine future-state prediction (a world model; most current mainstream WAMs output the next frame, similarly to video-generation models) with actions for policy planning.

Workflow: A decoupled architecture may be used, in which a policy model first outputs candidate actions and search-based rollouts in a world model optimize them. Alternatively, a sequential architecture may first use a world model to predict future states, then have a policy model output actions based on those predictions. The policy model and world model can also be combined into a single model with both functions.

Characteristics: This approach has risen rapidly in recent years, with mainstream methods generally fine-tuned from video-generation models. By forward-predicting state evolution, a world model allows the model to understand what deformations or displacements actions will cause in the physical environment, giving it the ability to ``plan'' in the real world. A model that interacts with an environment needs planning capability: when large models interact with the digital world, language models already serve a predictive and planning role; when embodied AI interacts with the physical world, world models serve that planning role.

\Needspace{5\baselineskip}
\hypertarget{ii.-the-actions-generated-by-embodied-ai-models}{%
\subsection{II. The ``Actions'' Generated by Embodied AI Models}\label{ii.-the-actions-generated-by-embodied-ai-models}}

\Needspace{5\baselineskip}
\hypertarget{what-exactly-are-actions}{%
\subsubsection{(1) What Exactly Are ``Actions''?}\label{what-exactly-are-actions}}

\Needspace{5\baselineskip}
Taking OpenVLA as an example, its actions use changes in the pose of the end effector (EEF). There are seven dimensions, representing the typical ``6-DoF end-effector motion + 1-DoF gripper'':

\[
a_t=[\Delta x,\Delta y,\Delta z,\Delta r_x,\Delta r_y,\Delta r_z,g].
\]

The OpenVLA authors explicitly state that the first three dimensions are relative changes in the end effector's XYZ position, the middle three are relative changes in its orientation, and the last controls the gripper. The specific orientation representation and units are defined by the corresponding dataset.

These are pose changes, not control commands for individual joints. An important purpose of this design is to decouple the task policy from the particular robot embodiment as far as possible. The same action, ``reach forward five centimeters and grasp the cup,'' can be executed by arms with different degrees of freedom and link lengths, without requiring the embodied large model to learn each robot's specific joint angles separately.

\Needspace{5\baselineskip}
While adopting a similar idea, \(\pi_0\) is relatively flexible in its total dimensionality, with a maximum of 32 dimensions. An example layout is:

\[
\begin{aligned}
\mathit{dim}_{0:5} &= \text{left arm joint angles},\\
\mathit{dim}_{6} &= \text{left gripper},\\
\mathit{dim}_{7:12} &= \text{right arm joint angles},\\
\mathit{dim}_{13} &= \text{right gripper}.
\end{aligned}
\]

\Needspace{5\baselineskip}
Mobile robots can additionally use:

\[
\mathit{dim}_{14:15}=[v_x,v_y].
\]

Other unused dimensions can be padding. For a 7-DoF single-arm robot such as Franka, the official documentation also allows the first seven dimensions to be used for joint actions and the eighth for the gripper.

The remaining dimensions accommodate different robot embodiments: some robots may use certain dimensions, while others use different ones, with unused dimensions set as padding. For some 7-DoF single-arm robots, the first seven dimensions may likewise be used for joint actions and the eighth for the gripper.

This is a very different design philosophy from OpenVLA's ``strictly seven action tokens.''

\Needspace{5\baselineskip}
As another example, LingBot-VA 2.0 has a maximum action dimensionality of 30, allocated as follows:

\Needspace{5\baselineskip}
For each hand:

\[
[\underbrace{\text{6-DoF root pose}}_{6},\underbrace{\mathit{gripper}}_{1}]
\]

\Needspace{5\baselineskip}
This can be abstracted as:

\[
[\mathit{pose}^{6D}_L,g_L,\mathit{pose}^{6D}_R,g_R].
\]

As with \(\pi_0\), the other 16 dimensions accommodate other embodiments and related requirements. Dimensions absent in a particular robot are masked out, making the design relatively flexible.

\Needspace{5\baselineskip}
\hypertarget{how-are-actions-converted-into-low-level-motor-commands}{%
\subsubsection{(2) How Are ``Actions'' Converted into Low-Level Motor Commands?}\label{how-are-actions-converted-into-low-level-motor-commands}}

As noted earlier, the ``actions'' output by embodied large models are not low-level commands, such as joint torques or currents, that can be sent directly to motor drivers. They are relatively general intermediate action representations, most typically changes in end-effector pose. Therefore, one or even several control layers often sit between the model's action generation and actual robot actuation. For example, when the model outputs an end-effector pose change, the robot first computes a target end-effector pose from its current position, then uses inverse kinematics (IK) to solve for the target angles of the joints. Under velocity control, end-effector velocity can instead be converted into joint velocities. The robot's own joint-position, velocity, impedance, or other controllers then produce the torque, current, and other commands required for actual motor execution.

For many robotic arms, the robot's URDF, link dimensions, joint constraints, Jacobian, and other parameters are usually known. Thus, converting general action representations into specific joint actions and motor commands does not require training an additional neural network on the spot; the robot's existing kinematic model and control algorithms can be used directly.

However, an important distinction is that the action-conversion module's not requiring training does not mean that the model necessarily needs no embodiment-specific adaptation training. For models supporting multiple robot embodiments, such as \(\pi_0\) and LingBot-VA 2.0, a unified action space only resolves inconsistencies in action dimensions and representations across robots. For example, the model can designate some dimensions for left-arm end-effector pose, others for right-arm end-effector pose, and others for joint positions, while a particular robot uses only the relevant dimensions. But a model's ability to ``accommodate'' a robot dimensionally does not mean that it inherently knows how to control that robot. A new robot unseen during training may have different camera mounting positions, arm workspaces, joint degrees of freedom, action scales, or control frequencies. Some degree of embodiment adaptation is therefore often required, such as fine-tuning or post-training on a small amount of demonstration data from that robot, or training a smaller action head/adapter.

Adaptation for real-robot deployment can therefore be divided into two levels.

The first is policy-level adaptation: enabling the embodied large model to understand how this robot should complete tasks. This level requires training when the embodiment, observation method, or action space differs substantially from those of the pretrained model. For example, a model facing a new robot not seen in pretraining can be fine-tuned on a small number of real-robot trajectories to adapt to its vision, proprioceptive states, and action distribution.

The second is control-level conversion: converting the model's end-effector poses, joint targets, and other action outputs into low-level control quantities that the robot can actually execute. This level typically does not require ad hoc training and relies more on traditional robotics control methods such as IK, whole-body control, QP, MPC, and impedance control.

Of course, in a few simple cases, the model may directly predict joint-space actions. For some 7-DoF robotic arms, for example, its outputs can be sent directly to the robot controller as joint-position or joint-increment targets. Even then, however, model outputs generally do not directly become motor currents; traditional control methods usually still need to compute actual joint torques from them.

\FloatBarrier
\Needspace{5\baselineskip}
\hypertarget{references-150}{%
\subsection{References}\label{references-150}}

\vspace{0.25em}
\begin{itemize}
\tightlist
\item
  Kim, M. J., Pertsch, K., Karamcheti, S., et al.~(2024). \href{https://arxiv.org/abs/2406.09246}{OpenVLA: An Open-Source Vision-Language-Action Model}. arXiv:2406.09246.
\item
  Black, K., Brown, N., Driess, D., et al.~(2024). \(\pi_0\): \href{https://arxiv.org/abs/2410.24164}{A Vision-Language-Action Flow Model for General Robot Control}. arXiv:2410.24164.
\item
  Zhang, Q., Li, L., Zhang, L., et al.~(2026). \href{https://arxiv.org/abs/2607.08639}{Native Video-Action Pretraining for Generalizable Robot Control}. arXiv:2607.08639. (LingBot-VA 2.0)
\item
  Physical Intelligence. (2025). \href{https://github.com/Physical-Intelligence/openpi}{OpenPI: Open-Source Models and Packages for Robotics}. GitHub.
\item
  Physical Intelligence. (2025). \href{https://github.com/Physical-Intelligence/openpi/blob/main/docs/norm_stats.md}{Action Normalization Statistics and Robot-Specific Action Spaces}. OpenPI Documentation.
\item
  Hugging Face LeRobot. (2026). \href{https://huggingface.co/docs/lerobot/main/lingbot_va}{LingBot-VA Documentation}. LeRobot Documentation.
\item
  Dyna Robotics. (2026). \href{https://www.dyna.co/dyna-2}{Dyna-2: A 1-Million-Hour Scaling Law for World-Action Models}. Dyna Robotics.
\item
  Generalist. (2026). \href{https://generalistai.com/blog/gen-1}{GEN-1: Scaling Embodied Foundation Models to Mastery}. Generalist AI.
\item
  Generalist. (2026). \href{https://generalistai.com/blog/gen-1.5}{GEN-1.5: Embodied Foundation Models Are One-Shot Learners}. Generalist AI.
\item
  Khatib, O. (1987). A Unified Approach for Motion and Force Control of Robot Manipulators: The Operational Space Formulation. \emph{IEEE Journal on Robotics and Automation}, 3(1), 43--53.
\item
  Khatib, O. (1993). The Operational Space Framework. \emph{JSME International Journal Series C}, 36(3), 277--287.
\item
  Khatib, O., Sentis, L., Park, J., \& Warren, J. (2004). \href{https://khatib.stanford.edu/publications/pdfs/Khatib_2004_IJHR.pdf}{Whole-Body Dynamic Behavior and Control of Human-Like Robots}. \emph{International Journal of Humanoid Robotics}, 1(1), 29--43.
\item
  Escande, A., Mansard, N., \& Wieber, P. B. (2014). Hierarchical Quadratic Programming: Fast Online Humanoid-Robot Motion Generation. \emph{The International Journal of Robotics Research}, 33(7), 1006--1028.
\item
  Redmon, J., Divvala, S., Girshick, R., \& Farhadi, A. (2016). \href{https://openaccess.thecvf.com/content_cvpr_2016/html/Redmon_You_Only_Look_CVPR_2016_paper.html}{You Only Look Once: Unified, Real-Time Object Detection}. \emph{CVPR}, 779--788.
\item
  Cadena, C., Carlone, L., Carrillo, H., et al.~(2016). \href{https://arxiv.org/abs/1606.05830}{Past, Present, and Future of Simultaneous Localization and Mapping: Toward the Robust-Perception Age}. \emph{IEEE Transactions on Robotics}, 32(6), 1309--1332.
\item
  Ichter, B., Brohan, A., Chebotar, Y., et al.~(2023). \href{https://proceedings.mlr.press/v205/ichter23a.html}{Do As I Can, Not As I Say: Grounding Language in Robotic Affordances}. \emph{Conference on Robot Learning}, PMLR 205:287--318.
\item
  Liang, J., Huang, W., Xia, F., et al.~(2022). \href{https://arxiv.org/abs/2209.07753}{Code as Policies: Language Model Programs for Embodied Control}. arXiv:2209.07753.
\item
  Driess, D., Xia, F., Sajjadi, M. S. M., et al.~(2023). \href{https://proceedings.mlr.press/v202/driess23a.html}{PaLM-E: An Embodied Multimodal Language Model}. \emph{ICML}, PMLR 202:8469--8488.
\item
  Zitkovich, B., Yu, T., Xu, S., et al.~(2023). \href{https://proceedings.mlr.press/v229/zitkovich23a.html}{RT-2: Vision-Language-Action Models Transfer Web Knowledge to Robotic Control}. \emph{Conference on Robot Learning}, PMLR 229:2165--2183.
\item
  Wang, S., Shi, J., Fu, Z., et al.~(2026). \href{https://arxiv.org/abs/2605.12090}{World Action Models: The Next Frontier in Embodied AI}. arXiv:2605.12090.
\item
  Doshi, R., Walke, H. R., Mees, O., Dasari, S., \& Levine, S. (2025). \href{https://proceedings.mlr.press/v270/doshi25a.html}{Scaling Cross-Embodied Learning: One Policy for Manipulation, Navigation, Locomotion and Aviation}. \emph{Conference on Robot Learning}, PMLR 270:496--512.
\item
  Open Robotics. (n.d.). \href{https://docs.ros.org/en/rolling/Tutorials/Intermediate/URDF/URDF-Main.html}{URDF --- ROS 2 Documentation}. ROS 2 Documentation.
\item
  Whitney, D. E. (1969). \href{https://ieeexplore.ieee.org/document/4081862/}{Resolved Motion Rate Control of Manipulators and Human Prostheses}. \emph{IEEE Transactions on Man-Machine Systems}, 10(2), 47--53.
\item
  Hogan, N. (1985). \href{https://doi.org/10.1115/1.3140702}{Impedance Control: An Approach to Manipulation, Part I---Theory}. \emph{Journal of Dynamic Systems, Measurement, and Control}, 107(1), 1--7.
\item
  Mayne, D. Q., Rawlings, J. B., Rao, C. V., \& Scokaert, P. O. M. (2000). \href{https://www.sciencedirect.com/science/article/pii/S0005109899002149}{Constrained Model Predictive Control: Stability and Optimality}. \emph{Automatica}, 36(6), 789--814.
\item
  Åström, K. J., \& Hägglund, T. (2001). \href{https://www.sciencedirect.com/science/article/pii/S0967066101000624}{The Future of PID Control}. \emph{Control Engineering Practice}, 9(11), 1163--1175.
\end{itemize}

\clearpage
